\documentclass[12 pt]{extarticle}

	\usepackage[french]{babel}
	\usepackage[utf8]{inputenc}  
	\usepackage[T1]{fontenc}
	\usepackage{amssymb}
	\usepackage[mathscr]{euscript}
	\usepackage{stmaryrd}
	\usepackage{amsmath}
	\usepackage{tikz}
	\usepackage[all,cmtip]{xy}
	\usepackage{amsthm}
	\usepackage{varioref}
	\usepackage{geometry}
	\geometry{a4paper}
	\usepackage{lmodern}
	\usepackage{hyperref}
	\usepackage{array}
	 \usepackage{fancyhdr}
\renewcommand{\theenumi}{\alph{enumi})}
	\pagestyle{fancy}
	\theoremstyle{plain}
	\fancyfoot[C]{} 
	\fancyhead[L]{Contrôle}
	\fancyhead[R]{28 février 2026}\geometry{
 a4paper,
 total={170mm,257mm},
 left=20mm,
 top=20mm,
 }
	
	
	\title{Corrigé}
	\date{}
	\begin{document}
	\subsection*{Question 1}
	\begin{enumerate}
\item	$17-12+13-18 = 5 + (-5) = 0$ 
 \item $5+4-8- 1+4 - 2 = 9 - 9 + 2 = 2$
 \end{enumerate}
 \subsection*{Question 2}
 \begin{enumerate}
 \item $\frac49\div\frac29 = \frac49 \times \frac92  = \frac{4\times 9}{2\times 9} = \frac42 = 2$. 
 \item $\frac23\div\frac54 = \frac23\times \frac45 = \frac{2\times 4}{3\times 5} = \frac8{15}$.  
 \end{enumerate}
 \subsection*{Question 3}
\begin{enumerate} \item $2\div \frac13 = 2 \times 3 = 6$
\item $\frac16\div 2 = \frac16\times \frac12 = \frac1{6\times 2} = \frac12$
\item $3\times \frac{-3}4= \frac{3\times (-3}4 = - \frac94$.
\end{enumerate}
\subsection*{Question 4}
\begin{enumerate}
\item $(-2)^3\div 9 = - 8\div 9 = -\frac89$.
\item $(-3)^2 = (-3)\times(-3)=+9$.
\end{enumerate}

\subsection*{Question 5}
\begin{enumerate}
\item $2~350~000= 2,35\times 10^6$.
\item $103~258~202 = 1,032~582~02 \times 10^8$. 
\item $954~316~372,135 = 9,543~163~721~35\times 10^8$.
\item $3~298,204 = 3,298~204 \times 10^3$. 
\end{enumerate}
\subsection*{Question 6}
\begin{enumerate}
\item Non car $91= 7\times 13$. 
\item Oui, car $\sqrt{97} \simeq 9,\cdots$ et $97$ est impair (il finit par $7$), il n'est pas multiple de $3$ (car $9+7=16$ non plus), ni de $5$ (il ne finit pas par $0$ ou $5$) ni de $7$ ($97\div 7 = 13$ reste $6$), et le nombre premier suivant serait $11>\sqrt{97}$.
\end{enumerate}

\subsection*{Question 7}
\begin{enumerate}
\item $250 = 2 \times 125 = 2\times 5 \times 25 = 2\times 5 \times 5\times 5 = 2\times 5^3$.
\item $108= 2\times 54 = 2^2\times 27 = 2^2\times 3 \times 9 = 2^2\times 3\times 3 \times 3 = 2^2\times 3^3$.
\end{enumerate}
\subsection*{Question 8}
\begin{enumerate}
\item Pour $a=-2$, $5a+1= 5\times (-2)+1 = -10+1 = -9$.
\item Pour $a=-2$, $7a-4 = 7\times(-2)-4 = -14-4=-18$.
\item Pour $a=-3$, $7a-4 = 7\times(-3)-4 = -21-4=-25$.
\end{enumerate}
\subsection*{Question 9}
\begin{enumerate}
\item $(-3ac)(-2a^2) = (-3)\times (-2)aa^2c = 6 a^3x$.
\item $(7ab)(-2a^3) = -14a^4b$.
\item $(-abc)\times(2bc) = -2ab^2c^2$
\end{enumerate}
\subsection*{Question 10}
\begin{enumerate}
\item $3a+4b+12b-19a = (3-19)a + (4+12)b = -16a
+16b$.
\item $-8a+b-b^2+a = -7a+b-b^2$.
\item $12ac+41a-32ca+12c-4a = 41a-4a+12ac-32ac+12c= 37a-20ac+12c$.
\item $7x^2+4x-3-12x+53+12x^2 - 1 = (7+12)x^2 + (4-12)x + 53-3-1 = 19x^2-8x+49$.
\end{enumerate}
\subsection*{Question 11}
\begin{enumerate}
\item $-(-x+7) = -(-x) - (+7) = +x - 7 $
\item $-(-3x-3) = +3x+3$.
\item $-(x+1) = -x -1$.
\end{enumerate}
\subsection*{Question 12}
\begin{enumerate}
\item $-3(x+5) - 4(x+2) = -3\times x + (-3) \times 5 + (-4)\times x+ (-4) \times 2 = -3x-15-4x-8 = -7x-23$.
\item $x+4(2-x)+2 = x+ 4\times 2 + 4\times (-x) + 2 
= x + 8 - 4x +2 = -3x+10$.
\item $(x+4)-5(x+1) = x+4 - (5x+5) = x+4-5x-5 = -4x-1$.
\end{enumerate}
\subsection*{Question 13}
\begin{enumerate}
\item \begin{align*}&7(x+4)-(3x+1)+x(x+2)\\&= 
7\times x + 7\times 4 - 3x - 1 + x \times x + x \times 2\\ 
&= 7x + 28- 3x-1 +x^2 + 2x \\
&= x^2 +7x+2x-3x + 28-1 \\
&= x^2 + 6x+27\end{align*}
\item \begin{align*}
-(x+1)x - x -1 &= -(x\times x + 1\times x) - x -1 \\
&= -(x^2 + x) - x -1 \\
&= -x^2-x-x-1 \\
&= -x^2-2x-1
\end{align*}
\item \begin{align*}
x+x(1-x) -2x(x+1) & = x+ x\times 1 +x \times (-x)
- (2x\times x + 2x \times 1)\\
&= x + x -x^2 -(2x^2 +2x) \\
&= 2x - x^2 - 2x^2 - 2x \\
&= -3x^2
\end{align*}
\end{enumerate}
\subsection*{Question 14}
\begin{enumerate}
\item $15ac+3ab+6a = 3a\times(5c) + 3a b + 3a \times 2 = 
3a(5c+b+2)$.
\item $3x+xy-5xz = x(3+y-5z)$.
\item $4xy+12x = 4x(y+3)$.
\item $10xy+20yz = 10y(x+2z)$.
\end{enumerate}
\subsection*{Question 15}
\begin{enumerate}
\item $(x+1)(x+2)+(x+2)(x+3) = 
(x+2)[(x+1)+(x+3)]
= (x+2)(2x+4)$
\item $2(x+1) + 3(x+1) + x(x+1) = [2+3+x](x+1) = (5+x)(x+1)$.
\item $(x+1)(x+3)+(x+3)(2x-1) =
(x+3)[(x+1)+(2x-1)] 
= (x+3)[3x] = 3x(x+3)$.
\end{enumerate}
\subsection*{Question 16}
\begin{enumerate}
\item \begin{eqnarray*}(-x-5)(12x+1) &=& (-x)\times(12x) + (-x)\times 1
+ (-5)\times(12x) + (-5) \times 1 \\
&=& -12x^2 -x -60x -5 \\
&=& -12 x^2-61 x -5\end{eqnarray*}
\item \begin{eqnarray*}
(2x-3y)(4x-2) &=& 2x\times 4x + 2x\times (-2) + (-3y) \times 4x + (-3y)\times(-2)\\
&=& 8x^2 - 4x -12xy +6y 
\end{eqnarray*}
\item \begin{eqnarray*}
(2a+3b)(-4a+6b) &=& 
(2a)\times(-4a) + (2a)\times(6b) + (3b)\times(-4a) + (3b)\times (6b) \\
&=& -8a^2 +12ab -12ab + 18b^2 \\
&=& -8a^2+18b^2
\end{eqnarray*}
\item \begin{eqnarray*}
(x+2)(2x+2) &=& 
x\times 2x + x \times 2 + 2 \times 2x + 2 \times 2 \\
&=& 2x^2 + 2x + 4x + 4 \\
&=& 2x^2+6x+4
\end{eqnarray*}
\end{enumerate}
\subsection*{Question 17}
\begin{enumerate}
\item 
\begin{eqnarray*}(-3xy-2y)(x+5)& = &
(-3xy)\times x + (-3xy) \times 5
+ (-2y) \times x 
+ (-2y) \times 5 \\&=& 
-3x^2y -15xy -2xy -10 y \\
&=& -3x^2y-17xy-10y\end{eqnarray*}
\item \begin{eqnarray*}
(-4x+1)(y-3) &=& -4x\times y + (-4x) \times (-3) + 1\times y + 1\times (-3) \\
&=& -4xy +12x+y-3
\end{eqnarray*}
\end{enumerate}
\subsection*{Question 18}
\begin{enumerate}
\item On écrit successivement 
\begin{align*}
4x + 7 &= 51 \\
4x + 7 - 7 &= 51- 7 \\
4x &= 44 \\ 
\frac{4x}4  &=\frac{44}4\\
x &= 11 
\end{align*}
On vérifie aussi que $4\times 11+ 7 = 44+7 = 51$, donc $11$ est l'unique solution.
\item On écrit successivement 
\begin{align*}
7x + 8 &= 43 \\
7x+8 -8 &= 43-8\\
7x &= 35 \\ 
\frac{7x}7  &=\frac{35}7\\
x &= 5
\end{align*}
On vérifie aussi que $7\times 5+8 = 35+8 = 43$ donc $5$ est l'unique solution.
\item On écrit successivement 
\begin{align*}
2x-3 &= 19 \\
2x-3+3 &= 19+3\\
2x &= 22 \\ 
\frac{2x}2  &=\frac{22}2\\
x &= 11
\end{align*}
On vérifie aussi que $2\times11-3 = 22-3=19$ donc $11$ est l'unique solution.
\item On écrit successivement 
\begin{align*}
5x-23 &= 77 \\
5x-23+23 &= 77+23\\
5x &= 100 \\ 
\frac{5x}5  &=\frac{100}5\\
x &= 20
\end{align*}
On vérifie aussi que $5\times 20 - 23 = 100 - 23 = 77$ donc $20$ est l'unique solution.
\end{enumerate}
\subsection*{Question 19}
\begin{enumerate}
\item On a $1~cm^3 = 0, 001~dm^3$ donc $50~cm^3 = 0,05~dm^3$
\item On a $1~dm^3=1~000~cm^3$ donc $40~dm^3=40~000~cm^3$.
\item On a $1~m^3 = 1~000~dm^3$ donc $0,3~m^3 = 300~dm^3$.
\end{enumerate}
\subsection*{Question 20}
\begin{enumerate}
\item Le volume est $(3~cm)\times (5~cm)\times(15~cm) = 
225~cm^3 = 225~mL$.
\item Le volume est $(3~cm)\times (5~cm)\times(4~cm) = 
60~cm^3 = 60~mL$.
\end{enumerate}
 	\end{document}
